\subsection{Actual rational functions and finite-boundary maximum modulus}\label{subsec:fc-line-rational-functions}

\begin{definition}[The actual rational algebra and its evaluations]\label{def:fc-line-rational-functions}
\lean{EnvelopeProperty.Berkovich.RelativeSpectrum.FCStandardAffinoid.denominatorPolynomial, EnvelopeProperty.Berkovich.RelativeSpectrum.FCStandardAffinoid.denominatorPolynomial_ne_zero, EnvelopeProperty.Berkovich.RelativeSpectrum.FCStandardAffinoid.denominatorPolynomial_val_ne_zero, EnvelopeProperty.Berkovich.RelativeSpectrum.FCStandardAffinoid.RationalFunctions, EnvelopeProperty.Berkovich.RelativeSpectrum.FCStandardAffinoid.rationalPoint, EnvelopeProperty.Berkovich.RelativeSpectrum.FCStandardAffinoid.rationalValue, EnvelopeProperty.Berkovich.RelativeSpectrum.FCStandardAffinoid.rationalValue_polynomial, EnvelopeProperty.Berkovich.RelativeSpectrum.FCStandardAffinoid.rationalValue_mk', EnvelopeProperty.Berkovich.RelativeSpectrum.FCStandardAffinoid.continuous_rationalValue, EnvelopeProperty.Berkovich.RelativeSpectrum.FCStandardAffinoid.rationalEvaluation, EnvelopeProperty.Berkovich.RelativeSpectrum.FCStandardAffinoid.rationalEvaluation_polynomial, EnvelopeProperty.Berkovich.RelativeSpectrum.FCStandardAffinoid.nnnorm_rationalEvaluation}
\leanok
\uses{def:fc-standard-affinoid}
Put $Q_U(T)=\prod_{j=1}^m(T-a_j)$ and
$R_U=F[T,Q_U^{-1}]=F[T,(T-a_1)^{-1},\ldots,(T-a_m)^{-1}]$.
Every point $x\in U$ defines its genuine multiplicative seminorm on $R_U$
and a ring homomorphism $R_U\to\mathcal H(x)$ extending the original
polynomial evaluation. These give the same value on every rational function.
\end{definition}
\begin{proof}
Every hole coordinate has value at least its positive hole radius on $U$.
Thus $|Q_U|_x$ is a finite product of positive values and is nonzero.
Localization extends the original valuation and the map to the original
completed residue field uniquely. Their values agree because both restrict
to the original polynomial seminorm and a localization seminorm is determined
by its restriction. In particular, on $P/Q_U^n$ both values are
$|P|_x/|Q_U|_x^n$. This also proves continuity of every rational-function
value on $U$, since its denominator is everywhere nonzero there.
The empty product is one, so the construction includes no holes.
\end{proof}

\begin{lemma}[Positive rational maxima are locally constant]\label{lem:fc-positive-rational-local-maximum}
\lean{EnvelopeProperty.Berkovich.RelativeSpectrum.FCStandardAffinoid.rationalRatio, EnvelopeProperty.Berkovich.RelativeSpectrum.FCStandardAffinoid.rationalRatio_eventually_eq_of_isLocalMax}
\leanok
\uses{def:fc-line-rational-functions}
Let $P,Q\in F[T]$. At an actual point where $|P|,|Q|$ are nonzero, a local
maximum of $|P|/|Q|$ is locally constant on an open neighborhood contained
in that same nonvanishing locus.
\end{lemma}
\begin{proof}
On the open locus where $PQ$ has nonzero value, the actual regular-section
unit logarithm is $\log|P|-\log|Q|=\log(|P|/|Q|)$.
Monotonicity of the real logarithm transfers the given local maximum to this
actual rational-unit logarithm. The proved all-point rational-unit local
maximum theorem gives an open neighborhood on which this same logarithm is
constant and both factors remain nonzero. Injectivity of the logarithm on
positive real numbers makes the original ratio constant there.
\end{proof}

\begin{theorem}[Finite-boundary maximum for actual rational functions]\label{thm:fc-rational-boundary-maximum}
\lean{EnvelopeProperty.Berkovich.RelativeSpectrum.FCStandardAffinoid.rationalRatio_exists_boundary_max, EnvelopeProperty.Berkovich.RelativeSpectrum.FCStandardAffinoid.rationalValue_exists_boundary_max}
\leanok
\uses{def:fc-line-rational-functions, lem:fc-positive-rational-local-maximum, lem:fc-line-topology}
For every $f\in R_U$ there is an actual $z\in\partial U$ such that
$|f|_x\le |f|_z$ for every actual $x\in U$.
\end{theorem}
\begin{proof}
Write $f=P/Q_U^n$. Its value is continuous on the compact connected set $U$,
so it has a maximum $M$. If $M=0$, every value is zero and the outer Gauss
boundary point realizes the maximum. Suppose $M>0$ and no boundary point
realizes it. The maximum locus is nonempty and closed. Every point of it
then lies in the topological interior of $U$. The maximum is positive, so
both $P$ and $Q_U^n$ have nonzero value there. Interior containment makes it
a local maximum on the actual affine line. The preceding lemma makes the
original ratio constant on a neighborhood, so the maximum locus is open
in $U$. Connectedness makes it all of $U$, contrary to the assumption that
the original outer boundary point is not in it. Thus a boundary point
realizes the maximum. This argument permits zero outer radius, coincident
Gauss boundaries, touching holes and the empty hole family.
\end{proof}
